% ======================================================================
%  sec12.tex —— 第 12 节 $\cO_X$-模层
%  底本：原书 pp.61–66（印刷页 = PDF 页）
%  原书标题：12. O_X-module sheaves
%  体例：\section + \subsection 与原文编号一致；定理类环境紧跟 \label{类型:编号}。
% ======================================================================

\section{$\cO_X$-模层}\label{sec:12}

% ---------------------------------------------------------------- p.61
\subsection{定义}\label{sec:12.1}

讨论交换环而不提它的模，是不完整的。

\begin{definition}\label{def:12.1.1}
设 $(X,\cO_X)$ 是环化空间。\emph{$\cO_X$-模预层}\footnote{另一名称：\emph{presheaf of modules for $\cO_X$}（$\cO_X$ 的模预层）。}是指一个对象 $\cM\in\PShv(X,\Ab)$，连同下述结构：
\begin{itemize}
  \item 对任意开子集 $U\subset X$，阿贝尔群 $\cM(U)$ 带有 $\cO_X(U)$-模结构，
\end{itemize}
使得
\begin{itemize}
  \item 对任意开子集 $U\subseteq V\subseteq X$，限制映射
  \[
  \cM(V)\to\cM(U)
  \]
  与同态
  \[
  \cO_X(V)\to\cO_X(U)
  \]
  的作用相容。
\end{itemize}
\end{definition}

$\cO_X$-模（\emph{层}）是指底层对象 $\cM\in\PShv(X,\Ab)$ 为层的 $\cO_X$-模预层。无歧义时，径称 $\cO_X$-模。

\begin{example}\label{ex:12.1.2}
结构层 $\cO_X$ 本身是 $\cO_X$-模（层）。
\end{example}

\begin{definition}\label{def:12.1.3}
设 $\cM$ 与 $\cN$ 是 $\cO_X$-模预层。从 $\cM$ 到 $\cN$ 的\emph{$\cO_X$-线性态射}，是指层间的态射 $f\colon\cM\to\cN$，使得：
\begin{itemize}
  \item 对任意开子集 $U\subset X$，映射 $f_U\colon\cM(U)\to\cN(U)$ 是 $\cO_X(U)$-线性的。
\end{itemize}
\end{definition}

由此得到范畴
\[
\omodP,
\]
其对象是 $\cO_X$-模预层，态射是 $\cO_X$-线性态射。再记
\[
\omod\subseteq\omodP
\]
为 $\cO_X$-模（层）构成的全子范畴。

\begin{yizhu}
$\omod$ 是阿贝尔范畴：核可以逐开集取，余核则须在逐开集取商之后再层化。Stacks 另给出一条判据——复形 $\mathcal F\to\mathcal G\to\mathcal H$ 在 $\mathcal G$ 处正合，当且仅当对每个点 $x\in X$，茎上的复形 $\mathcal F_x\to\mathcal G_x\to\mathcal H_x$ 正合。正合性因此是逐茎可检验的局部性质。
[来源：Stacks Project, Tag 01AG（Lemma 17.3.1），\url{https://stacks.math.columbia.edu/tag/01AG}（2026-11-17 访问）。]
\end{yizhu}

\begin{example}\label{ex:12.1.4}
$\cO_X$ 的理想层 $\cI$ 带有显然的 $\cO_X$-模结构。嵌入 $\cI\to\cO_X$ 是 $\cO_X$-线性态射。
\end{example}

% ---------------------------------------------------------------- p.61–62
\subsection{层化}\label{sec:12.2}

\begin{proposition}\label{prop:12.2.1}
嵌入函子 $\omod\to\omodP$ 有左伴随
\[
\omodP\to\omod,
\]
且与层化函子
\[
(-)^{\sharp}\colon\PShv(X,\Ab)\to\Shv(X,\Ab)
\]
相容。
\end{proposition}

\begin{sketch}
设 $\cM\in\omodP$，考虑其层化 $\cM^{\sharp}$，并视之为 $\Shv(X,\Ab)$ 中的对象。$\cO_X$ 在 $\cM^{\sharp}$ 上的作用由公式
\[
\cO_X\times\cM^{\sharp}\simeq\cO_X^{\sharp}\times\cM^{\sharp}\simeq(\cO_X\times\cM)^{\sharp}\to\cM^{\sharp}
\]
给出，其中
\begin{itemize}
  \item 第一个同构是因为 $\cO_X$ 是层；
  \item 第二个同构是因为层化与有限积交换（见注 3.1.5）；
  \item 第三个态射由 $\cO_X$ 在 $\cM$ 上的作用给出。
\end{itemize}
容易验证，上述公式定义了 $\cO_X$-模层 $\cM^{\sharp}$，且态射 $\phi\colon\cM\to\cM^{\sharp}$ 是 $\cO_X$-线性的。此外，对任意测试对象 $\cN\in\omod$，复合
\[
\Hom_{\omod}(\cM^{\sharp},\cN)\xrightarrow{-\circ\phi}\Hom_{\omodP}(\cM,\cN)
\]
是双射。这就给出所要的命题。
\end{sketch}

\begin{yizhu}
一般情形下，环预层 $\cO$ 本身也须一并层化。Stacks 的 Lemma 6.20.1 断言：$\cO^{\sharp}\times\cF^{\sharp}\to\cF^{\sharp}$ 使 $\cF^{\sharp}$ 成为 $\cO^{\sharp}$-模层，且层化 $(-)^{\sharp}$ 与「把 $\cO^{\sharp}$-模层看成 $\cO$-模预层」的遗忘函子互为伴随。本书的 $\cO_X$ 已是层，故只需层化模预层；第一个同构 $\cO_X\simeq\cO_X^{\sharp}$ 正是这一点的体现。
[来源：Stacks Project, Tag 0088（Section 6.20, Sheafification of presheaves of modules），\url{https://stacks.math.columbia.edu/tag/0088}（2026-11-17 访问）。]
\end{yizhu}

% ---------------------------------------------------------------- p.62
\subsection{前推与拉回}\label{sec:12.3}

\begin{construction}\label{con:12.3.1}
设 $\phi\colon(X,\cO_X)\to(Y,\cO_Y)$ 是环化空间之间的态射。回忆 $\Shv(Y,\CRing)$ 中有态射 $\alpha\colon\cO_Y\to\phi_*\cO_X$。

对任意 $\cM\in\omodP$，容易验证对象
\[
\phi_*(\cM)\in\PShv(Y,\Ab)
\]
带有 $\phi_*(\cO_X)$-模预层结构，结构映射为
\[
\phi_*(\cO_X)\times\phi_*(\cM)\simeq\phi_*(\cO_X\times\cM)\to\phi_*(\cM).
\]
沿 $\alpha$ 限制，$\phi_*(\cM)$ 又得到 $\cO_Y$-模预层结构。所得对象
\[
\phi_*(\cM)\in\cO_Y\text{-}\mathsf{mod}^{\PShv}
\]
称为 $\cO_X$-模预层 $\cM$ 沿 $\phi$ 的\emph{前推}（即\emph{正像}）。

这一构造可以升级为函子
\[
\phi_*\colon\omodP\to\cO_Y\text{-}\mathsf{mod}^{\PShv},
\]
它与阿贝尔层的前推函子相容。此外，它限制为函子
\[
\phi_*\colon\omod\to\cO_Y\text{-}\mathsf{mod}.
\]
\end{construction}

\begin{construction}\label{con:12.3.2}
对链
\[
(X,\cO_X)\xrightarrow{\phi}(Y,\cO_Y)\xrightarrow{\varphi}(Z,\cO_Z),
\]
复合
\[
\omodP\xrightarrow{\phi_*}\cO_Y\text{-}\mathsf{mod}^{\PShv}\xrightarrow{\varphi_*}\cO_Z\text{-}\mathsf{mod}^{\PShv}
\]
与 $(\varphi\circ\phi)_*$ 有典范等同。
\end{construction}

\begin{construction}\label{con:12.3.3}
设 $\phi\colon(X,\cO_X)\to(Y,\cO_Y)$ 是环化空间之间的态射。回忆 $\Shv(X,\CRing)$ 中有态射 $\beta\colon\phi^{-1}\cO_Y\to\cO_X$。特别地，$\PShv(X,\CRing)$ 中有态射
\[
\phi^{-1}_{\PShv}\cO_Y\to\phi^{-1}\cO_Y\to\cO_X.
\]

对任意 $\cN\in\cO_Y\text{-}\mathsf{mod}^{\PShv}$，容易验证对象
\[
\phi^{-1}_{\PShv}(\cN)\in\PShv(X,\Ab)
\]
带有 $\phi^{-1}_{\PShv}(\cO_Y)$-模结构\footnote{第一个同构见注 3.3.4。}，其结构映射为
\[
\phi^{-1}_{\PShv}(\cO_Y)\times\phi^{-1}_{\PShv}(\cN)\simeq\phi^{-1}_{\PShv}(\cO_Y\times\cN)\to\phi^{-1}_{\PShv}(\cN).
\]
\end{construction}

% ---------------------------------------------------------------- p.63
考虑 $X$ 上的 $\cO_X$-模预层
\[
U\mapsto\cO_X(U)\otimes_{(\phi^{-1}_{\PShv}\cO_Y)(U)}(\phi^{-1}_{\PShv}\cN)(U).
\]
记之为
\[
\phi^{*}_{\PShv}(\cN)\in\omodP,
\]
并称之为 $\cO_Y$-模预层 $\cN$ 沿 $\phi$ 的\emph{拉回}（即\emph{逆像}）。

这一构造可以升级为函子
\[
\phi^*_{\PShv}\colon\cO_Y\text{-}\mathsf{mod}^{\PShv}\to\omodP.
\]

\begin{proposition}\label{prop:12.3.4}
设 $\phi\colon(X,\cO_X)\to(Y,\cO_Y)$ 是环化空间之间的态射。函子
\[
\phi^*_{\PShv}\colon\cO_Y\text{-}\mathsf{mod}^{\PShv}\to\omodP
\]
典范左伴随于
\[
\phi_*\colon\omodP\to\cO_Y\text{-}\mathsf{mod}^{\PShv}.
\]
\end{proposition}

\begin{sketch}
设 $\cM\in\omodP$，$\cN\in\cO_Y\text{-}\mathsf{mod}^{\PShv}$。按定义，给出态射 $\phi^*_{\PShv}(\cN)\to\cM$，等同于给出一族与限制映射相容的 $\cO_X(U)$-线性映射
\[
\cO_X(U)\otimes_{(\phi^{-1}_{\PShv}\cO_Y)(U)}(\phi^{-1}_{\PShv}\cN)(U)\to\cM(U).
\]
后者又等同于给出一族与限制映射相容的 $(\phi^{-1}_{\PShv}\cO_Y)(U)$-线性映射
\[
(\phi^{-1}_{\PShv}\cN)(U)\to\cM(U),
\]
即一个 $\phi^{-1}_{\PShv}\cO_Y$-线性态射
\[
\phi^{-1}_{\PShv}\cN\to\cM,
\]
其中 $\phi^{-1}_{\PShv}(\cO_Y)$ 通过同态 $\phi^{-1}_{\PShv}(\cO_Y)\to\cO_X$ 作用在 $\cM$ 上。利用伴随 $(\phi^{-1}_{\PShv},\phi_*)$，上述数据等同于一个 $\cO_Y$-线性态射
\[
\cN\to\phi_*\cM,
\]
此即所求。
\end{sketch}

下面定义 $\cO$-模层的拉回。

\begin{construction}\label{con:12.3.5}
设 $\phi\colon(X,\cO_X)\to(Y,\cO_Y)$ 是环化空间之间的态射。把 $\phi^*$ 定义为下述复合：
\[
\phi^*\colon\cO_Y\text{-}\mathsf{mod}\xrightarrow{\subseteq}\cO_Y\text{-}\mathsf{mod}^{\PShv}\xrightarrow{\phi^*_{\PShv}}\omodP\xrightarrow{\sharp}\omod.
\]
对 $\cN\in\cO_Y\text{-}\mathsf{mod}$，称 $\phi^*\cN$ 为 $\cN$ 沿 $\phi$ 的\emph{拉回}（即\emph{逆像}）。

换言之，$\phi^*\cN$ 是预层
\[
U\mapsto\cO_X(U)\otimes_{(\phi^{-1}_{\PShv}\cO_Y)(U)}(\phi^{-1}_{\PShv}\cN)(U)
\]
的层化。
\end{construction}

\begin{exercise}\label{exc:12.3.6}
证明 $\phi^*\cN$ 可以等同于
\begin{equation}\tag{12.1}\label{eq:12.1}
U\mapsto\cO_X(U)\otimes_{(\phi^{-1}\cO_Y)(U)}(\phi^{-1}\cN)(U)
\end{equation}
的层化。
\end{exercise}

作为上述习题的推论，有：

\begin{yizhu}
(12.1) 正是 Stacks 采用的拉回定义：$f^{*}\cG = \cO_X\otimes_{f^{-1}\cO_Y}f^{-1}\cG$。同一处还证明 $\Hom_{\cO_X}(f^{*}\cG,\cF)\simeq\Hom_{\cO_Y}(\cG,f_*\cF)$，即 $f^{*}$ 左伴随于 $f_*$；本书的推论 \ref{cor:12.3.11} 与之逐字对应。两式只差用 $\phi^{-1}_{\PShv}$ 还是 $\phi^{-1}$，层化之后二者相同——这正是习题 \ref{exc:12.3.6} 的内容。
[来源：Stacks Project, Tag 0094（Section 6.26, Morphisms of ringed spaces and modules），\url{https://stacks.math.columbia.edu/tag/0094}（2026-11-17 访问）。]
\end{yizhu}

% ---------------------------------------------------------------- p.64
\begin{lemma}\label{lem:12.3.7}
设 $j\colon U\to X$ 是概形的开浸入。对 $\cM\in\omod$，有 $j^*\cM\simeq j^{-1}\cM$。
\end{lemma}

\begin{warning}\label{warn:12.3.8}
(12.1) 所给的对应一般不是层。换言之，要定义 $\phi^*$，必须在取张量积之后再层化。
\end{warning}

\begin{exercise}\label{exc:12.3.9}
设 $k$ 是域，考虑概形间的态射 $\phi\colon X\to Y$：
\[
\bigsqcup_I \Spec(k[t])\to\Spec(k),
\]
其中 $I$ 是无限集。设 $\cN\in\cO_Y\text{-}\mathsf{mod}\simeq k\text{-}\mathsf{mod}$ 是无限维对象。证明 (12.1) 不是层。
\end{exercise}

\begin{exercise}\label{exc:12.3.10}
设 $k$ 是域。设 $\phi\colon\A^2_k\to\A^1_k$ 是由 $k[x]\to k[x,y]$ 给出的投影映射，$i\colon\Spec(k)\to\A^1_k$ 是由 $k[x]\to k[x]/x\simeq k$ 给出的闭浸入。考虑 $\cN:=i_*\cO_{\Spec(k)}$。证明 (12.1) 不是层。
\end{exercise}

下述结果由命题 \ref{prop:12.2.1} 与命题 \ref{prop:12.3.4} 形式地推出。

\begin{corollary}\label{cor:12.3.11}
设 $\phi\colon(X,\cO_X)\to(Y,\cO_Y)$ 是环化空间之间的态射。函子
\[
\phi^*\colon\cO_Y\text{-}\mathsf{mod}\to\omod
\]
典范左伴随于
\[
\phi_*\colon\omod\to\cO_Y\text{-}\mathsf{mod}.
\]
\end{corollary}

\begin{construction}\label{con:12.3.12}
对链
\[
(X,\cO_X)\xrightarrow{\phi}(Y,\cO_Y)\xrightarrow{\varphi}(Z,\cO_Z),
\]
典范自然同构
\[
\varphi_*\circ\phi_*\simeq(\varphi\circ\phi)_*
\]
诱导典范自然同构
\[
\varphi^{\PShv}_*\circ\phi^{\PShv}_*\simeq(\varphi\circ\phi)^{\PShv}_*,
\qquad
\phi^*\circ\varphi^*\simeq(\varphi\circ\phi)^*.
\]
\end{construction}

% ---------------------------------------------------------------- p.64–65
\subsection{张量积与内 Hom}\label{sec:12.4}

\begin{construction}\label{con:12.4.1}
设 $(X,\cO_X)$ 是环化空间。对 $\cM,\cN\in\omod$，定义
\[
\cM\otimes_{\cO_X}\cN\in\omod
\]
为 $\cO_X$-模预层
\[
U\mapsto\cM(U)\otimes_{\cO_X(U)}\cN(U)
\]
的层化。
\end{construction}

\begin{example}\label{ex:12.4.2}
有 $\cM\otimes_{\cO_X}\cO_X\simeq\cO_X\otimes_{\cO_X}\cM\simeq\cM$。
\end{example}

\begin{lemma}\label{lem:12.4.3}
设 $x\in X$。有
\[
(\cM\otimes_{\cO_X}\cN)_x\simeq\cM_x\otimes_{\cO_{X,x}}\cN_x.
\]
\end{lemma}

\begin{proof}
由下述两件事合起来即得：层化不改变茎；滤过余极限与张量积交换。
\end{proof}

\begin{yizhu}
要点有二：张量积先逐开集取张量、再层化；层化不改茎。两件事合起来，才使茎上的张量积等于逐开集张量的茎。Stacks 把这条计算单列为引理，理由是张量积层的局部性质几乎都由它推出，例如可逆层的纤维维数。
[来源：Stacks Project, Tag 01CA（Section 17.16, Tensor product），\url{https://stacks.math.columbia.edu/tag/01CA}（2026-11-17 访问）。]
\end{yizhu}

% ---------------------------------------------------------------- p.65
\begin{remark}\label{rem:12.4.4}
事实上，可以在 $\omod$ 上定义\emph{对称幺半结构}，使张量积函子由上述构造给出。
\end{remark}

\begin{construction}\label{con:12.4.5}
设 $(X,\cO_X)$ 是环化空间。对 $\cM,\cN\in\omod$，容易验证
\[
U\mapsto\uHom_{\cO_X|_U}(\cM|_U,\cN|_U)
\]
是 $\cO_X$-模层。这里截面 $f\in\cO_X(U)$ 通过前复合 $\cM|_U\xrightarrow{f}\cM|_U$ 作用，或等价地通过后复合 $\cN|_U\xrightarrow{f}\cN|_U$ 作用。

所得对象记为
\[
\uHom_{\cO_X}(\cM,\cN)\in\omod.
\]
\end{construction}

\begin{example}\label{ex:12.4.6}
有 $\uHom_{\cO_X}(\cO_X,\cN)\simeq\cN$。
\end{example}

\begin{construction}\label{con:12.4.7}
设 $x\in X$。对 $x$ 的任意开邻域 $U\subseteq X$，有典范映射
\[
\uHom_{\cO_X}(\cM,\cN)(U)=\Hom_{\cO_X|_U}(\cM|_U,\cN|_U)\to\Hom_{\cO_{X,x}}(\cM_x,\cN_x).
\]
这诱导态射
\begin{equation}\tag{12.2}\label{eq:12.2}
\uHom_{\cO_X}(\cM,\cN)_x\to\Hom_{\cO_{X,x}}(\cM_x,\cN_x).
\end{equation}
\end{construction}

\begin{warning}\label{warn:12.4.8}
典范态射 (12.2) 一般不是同构。例如取 $\cM=\bigoplus_I\cO_X$，该映射为
\[
\Bigl(\prod_I\cN\Bigr)_x\to\prod_I(\cN_x).
\]
\end{warning}

\begin{yizhu}
(12.2) 失效的根源在无穷积：茎是滤过余极限，只与有限积交换。取 $\cM=\bigoplus_I\cO_X$ 时，$\uHom_{\cO_X}(\cM,\cN)\simeq\prod_I\cN$，其茎为 $\bigl(\prod_I\cN\bigr)_x$；而 $\Hom_{\cO_{X,x}}(\cM_x,\cN_x)\simeq\prod_I\cN_x$。$I$ 无限时二者一般不同，Stacks 对同一现象的说法是 "rarely an isomorphism"。
[来源：Stacks Project, Tag 01CM（Section 17.22, Internal Hom），\url{https://stacks.math.columbia.edu/tag/01CM}（2026-11-17 访问）。]
\end{yizhu}

\begin{construction}\label{con:12.4.9}
设 $(X,\cO_X)$ 是环化空间。对 $\cM_i\in\omod$（$i=1,2,3$），有典范态射
\[
\uHom_{\cO_X}(\cM_1,\cM_2)\otimes_{\cO_X}\uHom_{\cO_X}(\cM_2,\cM_3)\to\uHom_{\cO_X}(\cM_1,\cM_3).
\]
取 $\cM_1=\cO_X$、$\cM_2=\cM$、$\cM_3=\cN$，得到关于 $\cN$ 函子性的态射
\[
\cM\otimes_{\cO_X}\uHom_{\cO_X}(\cM,\cN)\to\cN.
\]
换言之，有自然变换
\begin{equation}\tag{12.3}\label{eq:12.3}
\cM\otimes_{\cO_X}\uHom_{\cO_X}(\cM,-)\to\id.
\end{equation}
\end{construction}

下述结果的证明与交换代数中的对应结论相同。

\begin{proposition}\label{prop:12.4.10}
对任意 $\cM\in\omod$，自然变换 (12.3) 给出伴随对
\[
\cM\otimes_{\cO_X}-\colon\omod\rightleftarrows\omod\colon\uHom_{\cO_X}(\cM,-).
\]
\end{proposition}

\begin{yizhu}
伴随的另一半是层版张量–内 Hom 同构
\[
\uHom_{\cO_X}(\cM\otimes_{\cO_X}\cN,\cP)\simeq\uHom_{\cO_X}(\cM,\uHom_{\cO_X}(\cN,\cP)),
\]
Stacks 的 Lemma 17.22.1 给出它；把恒等态射 $\uHom_{\cO_X}(\cM,\cN)\to\uHom_{\cO_X}(\cM,\cN)$ 沿该同构搬到左边，即得 (12.3) 的求值态射。命题 \ref{prop:12.4.10} 只是把这句同构改写成伴随对的说法。
[来源：Stacks Project, Tag 01CM（Lemma 17.22.1），\url{https://stacks.math.columbia.edu/tag/01CM}（2026-11-17 访问）。]
\end{yizhu}

\begin{corollary}\label{cor:12.4.11}
对任意 $\cM\in\omod$，有：
\begin{itemize}
  \item 函子 $\cM\otimes_{\cO_X}-$ 与余极限交换。特别地，它右正合。
  \item 函子 $\uHom_{\cO_X}(\cM,-)$ 与极限交换。特别地，它左正合。
\end{itemize}
\end{corollary}

事实上，下述结果也不难证明：

% ---------------------------------------------------------------- p.66
\begin{lemma}\label{lem:12.4.12}
对任意 $\cN\in\omod$，函子
\[
\uHom_{\cO_X}(-,\cN)\colon\omod^{\op}\to\omod
\]
与极限交换。特别地，它左正合。
\end{lemma}
